\chapter{Modello Relazionale}
\section{Relazione matematica}
Siano $D_1,D_2,...,D_n$ n insiemi:\\
Il \textbf{prodotto cartesiano} $D_1 \times D_2 \times ... \times D_n$ è l'insieme di tutte le n-uple ordinate $(d_1,d_2,..,d_n)$ tali che $d_1 \in D_1, d_2 \in D_2 ,..., d_n \in D_n$. 
Una relazione matematica sugli insiemi $D_1,D_2,...,D_n$ è un \textbf{sottoinsieme del prodotto cartesiano} $D_1 \times D_2 \times ... \times D_n$.\\
Se R è una relaziome matematica sugli insiemi $D_1,D_2,...,D_n$ tali insiemi sono detti domini, tale relazione su $n$ domini dice che ha grado $n$. il numero di n-uple (righe) di una relazione è la cardinalità della relazione stessa.\\
Esempio:
\begin{itemize}

    \item $D_1 = \{\text{Conti}, \text{Totti}, \text{Maradona}\}$

    \item $D_2 = \{\text{Genoa}, \text{Roma}\}$

    \item prodotto cartesiano $D_1 \times D_2$
    
    \[
    \begin{array}{|l|l|}
    \hline
    \text{Conti}    & \text{Genoa} \\
    \text{Conti}    & \text{Roma}  \\
    \text{Totti}    & \text{Genoa} \\
    \text{Totti}    & \text{Roma}  \\
    \text{Maradona} & \text{Genoa} \\
    \text{Maradona} & \text{Roma}  \\
    \hline
    \end{array}
    \]

    \item una relazione:
    $\text{HaGiocato} \subseteq D_1 \times D_2$

    \[
    \begin{array}{|l|l|}
    \hline
    \text{Conti} & \text{Roma}  \\
    \text{Conti} & \text{Genoa} \\
    \text{Totti} & \text{Roma}  \\
    \hline
    \end{array}
    \]

\end{itemize}

\noindent La relazione matematica è quindi un insieme di ennuple o tuple ordinate su $D_1,D_2,...,D_n$, dove ogni tupla è una sequenza di valori della forma $(d_1,...,d_n)$.
Ciascuna tupla è ordinata nel senso che il valore del primo elemento della tupla proviene dal primo domonio, ma non è definito un ordinamento tra le tuple e sono tutte distinte.
